\subsection{THETA3-C28: the actual Schoen fibration cannot produce the exceptional cycle by base change alone}\label{C28}

For the particular target $A$ of C25, the fibration
construction has a stronger limitation than the conditional
test in C27.  Every map to $A$ from a surface obtained by
base change of the chosen Schoen curve fibration factors
through a curve.  This statement includes all maps on
those surfaces, not only the previously written Albanese
quotients.  It follows from the explicit generic fibre
and does not assume that an unspecified fixed part is CM.

\paragraph{The connected generic fibre.}
Use the Fermat projection of Schoen, printed p.~133,
\[
     f:Y\longrightarrow C_5,\qquad
             (y_{ij})\longmapsto(y_{01}:y_{03}:y_{13}),
\]
on the smooth ball quotient $Y$.  Its composition with
the degree-five quotient $C_5\to C_0$ is the map used
in C15--C16 and C25.  On the chart $y_{03}=1$, write
\[
 u=y_{13},\quad s=u^5,\quad v=y_{23},\quad
 b=y_{02},\quad c=y_{12}.
\]
The original defining equations give
\[
 y_{01}^5=1-s,\qquad
 v^5=t,\qquad b^5=1-t,\qquad c^5=s-t,
                  \qquad t=v^5.                 \tag{C28a}
\]
Consequently the geometric generic fibre is the smooth
projective normalization $C_s$ of the last three equations,
where $s\ne0,1$ is transcendental over $\C$.
The three classes $t,1-t,s-t$ are independent modulo fifth
powers in $\overline{\C(C_5)}(t)^\times$:
their valuations at $0,1,s$ show separately that each
exponent in a fifth-power relation is divisible by five.
Thus the function field has degree $5^3$ over the
$t$-line, with deck group
\[
                         G=(\Z/5\Z)^3.
\]
This also proves geometric connectedness of the generic
fibre.  Resolution of the original surface does not change
its function field or this normalized generic fibre.
Since the base curve is normal, the Stein factor of $f$
is therefore $C_5$ itself.

At $0,1,s,\infty$ the inertia order is five, and there
is no other ramification.  At infinity the three radicals
have the same simple pole up to units, so its inertia is
the diagonal cyclic subgroup, rather than a group of
order $125$.  Riemann--Hurwitz gives
\[
 2g(C_s)-2=125(-2)+4\cdot125(1-1/5)=150,
                     \qquad g(C_s)=76.          \tag{C28b}
\]

\paragraph{All rational character factors have dimension at most four.}
A nontrivial deck character is specified by
\[
 (a_0,a_s,a_1,a_\infty)\in\mathbf F_5^4,
       \qquad a_0+a_s+a_1+a_\infty=0.
\]
The last entry records its value on infinity inertia.
Let $r$ be the number of nonzero entries.  The cyclic
quotient $C_\chi=C_s/\ker\chi$ has degree five over
$\mathbf P^1$ and is branched at exactly these $r$
points, each with inertia five.  Hence
\[
                      g(C_\chi)=2(r-2).          \tag{C28c}
\]
Here $r\in\{2,3,4\}$; a single nonzero entry is
incompatible with the displayed sum condition.

The rational character orbits consist of
$\chi,\chi^2,\chi^3,\chi^4$.  Pullback from $C_\chi$
identifies its first cohomology with the sum
of those four eigenspaces in the first cohomology of $C_s$.
Here one may use \'etale cohomology with coefficients
$\Q_\ell(\zeta_5)$, for a fixed $\ell\ne5$, over
the algebraically closed generic field.  The orbit
idempotents themselves belong to $\Q[G]$ and act by
algebraic correspondences.  Thus no embedding of an
extension of $\C(C_5)$ into $\C$ fixing the constants
is being assumed.
Indeed pullback and trace multiply by $\deg(C_s/C_\chi)$,
and averaging over $\ker\chi$ is the projector onto
its invariants.  The trivial-character part is zero,
since $C_s/G=\mathbf P^1$.
Distinct rational orbits give disjoint summands, and
together they exhaust first cohomology.  The corresponding
quotient maps and pullbacks therefore give an isogeny
\[
             J(C_s)\ \sim\ \prod_{[\chi]} J(C_\chi),
                                                        \tag{C28d}
\]
omitting genus-zero factors.  One can see the isogeny
directly on first cohomology: the product map has an
isomorphism there, and therefore has finite kernel and
full-dimensional image.

For completeness the exact character counts are
\[
\begin{array}{c|ccc}
 r&2&3&4\\\hline
 \text{characters}&24&48&52\\
 \text{rational orbits}&6&12&13\\
 \dim J(C_\chi)&0&2&4 .
\end{array}                                             \tag{C28e}
\]
For $r=2$, choose the two positions and one nonzero
entry, giving $6\cdot4=24$.  For $r=3$, choose the
zero position and two nonzero entries whose sum is
nonzero, giving $4(4^2-4)=48$.  Subtracting these and
the trivial character from $5^3$ leaves $52$.
Division by four gives the orbit counts.  The dimensions
sum to $12\cdot2+13\cdot4=76$, as required.

Every positive-dimensional factor in (C28d) has a
unital action by $E_5=\Q(\zeta_5)$.  Its order-five
deck automorphism has no invariant first cohomology,
so its minimal polynomial on that cohomology is
$1+T+T^2+T^3+T^4$; this produces the stated field action.
In particular the twelve dimension-two factors are
CM abelian surfaces.  No assertion about finite
monodromy of the thirteen dimension-four factors is
needed for the argument below.

\paragraph{The C25 target cannot occur in a fibre Jacobian.}
The target $A$ is geometrically simple of dimension four,
with $\operatorname{End}^0(A)=\Q(i)$ by C25.
These properties persist over algebraically closed
extensions of the constant field.  A precise source
for the invariance of abelian subvarieties and
homomorphisms under regular field extension is Milne,
\emph{Abelian Varieties}, version~2.00, Corollary~16.4(a)--(b),
printed p.~68; passage to an arbitrary algebraically
closed extension follows by finite presentation.

If there were a nonzero homomorphism $J(C_s)\to A$
over the algebraic closure of the base function field,
(C28d) would give a nonzero homomorphism from one of
its factors to $A$.  Simplicity of $A$ makes that
homomorphism surjective.  A dimension-two factor
cannot surject.  A dimension-four factor would be
isogenous to $A$, and its $E_5$ action would induce
a unital embedding
\[
                  \Q(\zeta_5)\hookrightarrow\Q(i),
\]
which is impossible by the field degrees.  We have proved
\[
                         \Hom(J(C_s),A)=0.       \tag{C28f}
\]
This argument does not presume that a dimension-four
factor is simple; a surjection between equal-dimensional
abelian varieties is already an isogeny.

\paragraph{Descent of every map on every such base change.}
Let $D\to C_5$ be any dominant morphism of smooth
projective connected curves, and let $X_D$ be a smooth
projective model of the main component of
$Y\times_{C_5}D$, with its induced map to $D$.
Geometric connectedness of the generic fibre gives
a unique main component.  By (C28f), every morphism
$h:X_D\to A$ is constant on the geometric generic
fibre: a map from a smooth projective curve to an
abelian variety, after choosing a point and translating,
is induced by its Jacobian homomorphism.

The resulting constant point is rational over
$\C(D)$.  To check this, extend scalars to an algebraic
closure of that field; each field automorphism fixes
the original map and hence fixes its unique constant
value.  Thus it descends to a rational map $D\dashrightarrow A$.
Properness of $A$ extends this map across every point
of the smooth curve $D$ by the valuative criterion.
The composite $X_D\to D\to A$ equals $h$ on the
generic fibre and hence on a dense open subset;
separatedness of $A$ makes them equal everywhere.
Therefore every $h$ factors through $D$.

For the original base $C_0$, form the normalization
of $C_5\times_{C_0}D_0$ and apply the preceding
argument to each connected component.  When
$D_0\to C_0$ is connected finite \'etale, the
degree-five Galois cover $C_5\to C_0$ pulls back
either to a connected degree-five cover or to five
copies of $D_0$, since its monodromy subgroup has
order either five or one.  These are exactly the
possible main components of the surface base change.
In every case the conclusion is factorization through
the corresponding Stein base curve.

It follows that the full quartic $F_{X_D}$ of C27
is zero for every surface in this single-fibration
base-change tower: the image of each map has dimension
at most one, so the pullback of every degree-four
class integrates to zero.  This includes the C25
cover and every further cover obtained by changing
only this base curve.  A map-generated exceptional
cycle on the specified $A$ must therefore use a
surface cover that genuinely changes the generic
fibre, or leave this map construction in favour of
a higher correspondence.  The proof does not rule
out either route, and does not claim the construction
of the exceptional cycle. \hfill$\square$
